% GPT-6 Astra Pro; effort not exposed; /root/replacement; 22 November 2026.
\section{HD-F1-R14: all couplings of cohomology filtered by semihomogeneous RM bundles retain a trace kernel}
\label{sec:followup-rm-filtered-semihomogeneous}

Let $X,\theta,F$ be the genuine real-quadratic fourfold of
R07. Let $E$ be an actual bounded perfect object with full
character
\[
 \operatorname{ch}(E)=N e^{-B}(1-qA_1)(1-qA_2),
 \qquad N\ne0,\quad q>0,
 \quad A_i=\theta_i^2/2.
 \tag{F1R14.1}
\]
Suppose each nonzero standard cohomology sheaf $H^j(E)$
admits a finite filtration by simple semihomogeneous vector
bundles whose rational slopes belong to $F\theta$.
Then every lift under every coherent faithful finite point
action satisfies
\[
                       \dim\ker\sigma^2_{\widetilde E}\ge2.
 \tag{F1R14.2}
\]
In particular no such lift is weakly semiregular.
There is no restriction on the number of slopes or on the
cohomological amplitude. All actual Postnikov couplings are
included. Each cohomology sheaf in this statement is a
vector bundle, since it is an extension of vector bundles;
arbitrary torsion or non-semihomogeneous cohomology is not
included in the assertion.

The proof uses the actual filtered Hochschild action and
the actual nullhomotopy on a Postnikov model. We first
write the adjacent case's common cancellation map, which
also explains why two independent Yoneda image containments
alone would not suffice for the conclusion.

\paragraph{1. A common cancellation map for the adjacent triangle.}
For the actual triangle
\[
 A\longrightarrow E\longrightarrow B_0:=T[-1]
      \xrightarrow{\delta_0}A[1],\qquad
                  \delta\in\operatorname{Ext}^2(T,A),
 \tag{F1R14.3}
\]
the complexes of filtered and ordinary derived endomorphisms
fit into the exact triangle
\[
 \operatorname{RHom}(A,B_0)[-1]\longrightarrow
 \operatorname{REnd}_{\mathrm{fil}}(E)\longrightarrow
 \operatorname{REnd}(E)\longrightarrow
                         \operatorname{RHom}(A,B_0).
 \tag{F1R14.4}
\]
Here the filtered object is the actual arrow $A\to E$.
One direct construction of (F1R14.4) is the homotopy fibre
\[
 \operatorname{REnd}_{\mathrm{fil}}(E)
 =\operatorname{fib}\left(
 \operatorname{REnd}(A)\oplus\operatorname{REnd}(E)
 \xrightarrow{(u,v)\mapsto i u-v i}
                         \operatorname{RHom}(A,E)\right).
\]
Taking the fibre of its forgetful map to
$\operatorname{REnd}(E)$ gives
$\operatorname{fib}(\operatorname{RHom}(A,A)
\to\operatorname{RHom}(A,E))
=\operatorname{RHom}(A,B_0)[-1]$, proving the displayed
triangle with its actual forgetful morphism.

The characteristic evaluation of a Hochschild class
$a\in HH^2(X)$ has a lift to degree two of this filtered
endomorphism complex, with endpoint values
$\operatorname{ev}_A(a)$ and $\operatorname{ev}_{B_0}(a)$.
To specify the required coherence, use the actual diagonal
kernel representation of Hochschild cohomology. In
Buchweitz--Flenner, Proposition 5.2.2 and its proof,
$a$ is represented on a projective diagonal resolution
$\overline P$ by an actual degree-two map
$\overline a:\overline P\to\overline P[2]$, and evaluation
on a complex $M$ is the roof
\[
 M\xleftarrow[\simeq]{\,1\otimes\overline\mu\,}
 M\otimes\overline P
 \xrightarrow{\,1\otimes\overline a\,}
 (M\otimes\overline P)[2]
 \xrightarrow[\simeq]{\,1\otimes\overline\mu\,}M[2].
 \tag{F1R14.5}
\]
The same resolution and maps apply simultaneously to every
term and arrow of the actual filtration. The augmentation
is a quasiisomorphism on every filtered piece, so the roof
is invertible in the filtered derived category as well.
This constructs the filtered lift; it uses the resolution
maps, not merely the underlying graded-centre property.

If $\operatorname{ev}_E(a)=0$, the long exact sequence of
(F1R14.4) places this filtered class in the boundary image of
\[
 H^1\operatorname{RHom}(A,B_0)
                         =\operatorname{Hom}(A,T).
\]
Thus there is one actual morphism $f:A\to T$ such that
\[
          \operatorname{ev}_A(a)=\delta f,\qquad
          \operatorname{ev}_T(a)=f[2]\delta.
 \tag{F1R14.6}
\]
The signs and the common $f$ can be checked in the actual
cone model. Write its differential as an upper triangular
matrix with off-diagonal entry $\delta_0$ of degree one.
A degree-one lower-left map represents $f:A\to B_0[1]=T$.
Its endomorphism differential is $d h+h d$; the two diagonal
entries are $\delta_0[1]f$ and $f[1]\delta_0$. After
shifting the second endpoint, these are exactly (F1R14.6).
This refines R10's two separate image containments by an
explicit morphism of the filtered endomorphism complexes.

\paragraph{2. The actual slope blocks and the even cross-Ext indices.}
Write a slope as $\ell_s=\theta_{a_s}$, $a_s\in F$,
so its two components are
$\sigma_1(a_s)\theta_1+\sigma_2(a_s)\theta_2$.
For distinct slopes $s,t$, put $d=a_t-a_s\ne0$ and
\[
       \nu(s,t)=2\#\{i\in\{1,2\}:\sigma_i(d)<0\}
                         \in\{0,2,4\}.
 \tag{F1R14.7}
\]
For simple semihomogeneous bundles $Q_s,Q_t$ at those
slopes there are exact vanishings
\[
            \operatorname{Ext}^k_X(Q_s,Q_t)=0
                         \quad(k\ne\nu(s,t)).
 \tag{F1R14.8}
\]
Here is the actual isogeny and index calculation. Mukai's
isogeny-line construction, as used in R09, realizes each
$Q$ as $f_*L$ for an isogeny and a line bundle. The actual
pullback is the direct sum of the translates of $L$ by
the kernel of $f$. A connected component of the fibre
product of the finitely many such isogenies, with its
origin fixed, supplies a common isogeny $p:Y\to X$.
Every $p^*Q$ is consequently a direct sum of line bundles
with common first Chern class $p^*\ell_Q$. Thus
$p^*(Q_s^\vee\otimes Q_t)$ is a direct sum of line bundles
with first Chern class $p^*\theta_d$.

For every such line bundle $L$ and every homogeneous twist,
Riemann--Roch gives the exact index polynomial
\[
 \chi_Y\bigl(L\otimes p^*\mathcal O_X(\theta)^{\otimes m}\bigr)
       =\deg(p)D\,(m+\sigma_1(d))^2(m+\sigma_2(d))^2,
       \qquad D=\int_X\theta^4/24.
 \tag{F1R14.9}
\]
Both embeddings of $d$ are nonzero. Mumford's Vanishing
Theorem and Index Theorem therefore put all cohomology
in precisely the degree (F1R14.7): the index is the number
of positive roots, counted with multiplicity, of (F1R14.9).
The original author's proof is in \emph{Abelian Varieties},
\S16, Vanishing Theorem and Index Theorem; in the inspected
scan the respective statements/proofs occupy printed
pp.~150--153 and 155--162, with the complex Hermitian
signature corollary on pp.~162--163.
Finally the actual unit and normalized trace
\[
 \mathcal O_X\longrightarrow p_*\mathcal O_Y
       \xrightarrow{\operatorname{Tr}_p/\deg p}\mathcal O_X
\]
split the pullback map on each cohomology group of
$Q_s^\vee\otimes Q_t$. The vanishing on $Y$ hence proves
(F1R14.8) on the original $X$.

Let $\mathcal C_s$ be the collection of actual vector bundles
with finite filtrations by simple semihomogeneous bundles
of slope $\ell_s$. Applying the long exact Ext sequences
to such filtrations proves (F1R14.8) for all
$F_s\in\mathcal C_s$ and $F_t\in\mathcal C_t$, still with
the same single possible index. In particular cross-slope
$\operatorname{Ext}^1$ vanishes. Every cohomology bundle
in the statement therefore has an actual decomposition
\[
                  H^j(E)=\bigoplus_s Q_{j,s},
          \qquad Q_{j,s}\in\mathcal C_s,
          \qquad\operatorname{ch}(Q_{j,s})=r_{j,s}e^{\ell_s}.
 \tag{F1R14.10}
\]
To see the decomposition without asserting a categorical
separation, induct on the given filtration length. Once
the earlier filtered bundle is decomposed into its slope
blocks, an extension by a simple quotient of slope $s$
has zero component in every $\operatorname{Ext}^1$
against a different block. Splitting those actual zero
extension classes gives the displayed direct sum, and
the remaining extension stays in $\mathcal C_s$.

The consequence needed below is particularly explicit:
for $s\ne t$, every product
\[
 \operatorname{Ext}^2(F_t,F_s)\circ
                         \operatorname{Hom}(F_s,F_t)
\]
is zero, as is the reverse product. Indeed a nonzero
Hom term forces $\nu(s,t)=0$, while
$\nu(t,s)=4-\nu(s,t)=4$, so the displayed Ext group is
zero. If the Hom group is zero, the product is already
zero. No assumption on the actual coupling is involved.

\paragraph{3. The full Postnikov model and its lower-triangular homotopy.}
Put $H_j=H^j(E)$ and $F_j=H_j[-j]$. The actual canonical
triangles
\[
 \tau^{\le j-1}E\longrightarrow\tau^{\le j}E
                  \longrightarrow H_j[-j]
                  \longrightarrow(\tau^{\le j-1}E)[1]
 \tag{F1R14.11}
\]
construct a finite twisted-complex model with graded
object $\bigoplus_jF_j$ and upper triangular differential
\[
                d=d_0+\Delta,\qquad
       \Delta_{ij}\in
       \mathcal A^{0,j-i+1}(\mathcal Hom(H_j,H_i))
                         \quad(i<j).
 \tag{F1R14.12}
\]
All entries here are the actual Postnikov morphisms.
Explicitly, use the Dolbeault dg category of vector bundles,
whose morphism complexes compute their Ext groups. Starting
with the lowest nonzero cohomology bundle, represent the
connecting map in (F1R14.11) by a closed degree-one column
in the already constructed model and take its actual cone.
Induction gives (F1R14.12) and $d^2=0$. The usual long
exact Hom sequences of these cones identify the cohomology
of their matrix morphism complexes with the self-Ext of
the original objects. Thus this construction is a model
of the original $E$, not a new object with just its character.

Apply (F1R14.5) to the complete actual filtration. It gives
a degree-two upper triangular cocycle $z_a$ in the filtered
endomorphism complex, whose diagonal entries represent
$\operatorname{ev}_{F_j}(a)$. One can compute the filtered
endomorphism complex by these upper triangular matrices:
for two steps this is exactly the homotopy fibre in
(F1R14.4), and iteration of that fibre description proves
the assertion for finitely many steps. The same diagonal
resolution and filtered quasiisomorphisms in (F1R14.5)
therefore supply the stated upper triangular representative.

Now let $a\in K_E=\ker c^2_{\operatorname{ch}(E)}$ and
assume its ordinary characteristic evaluation vanishes.
In the actual ordinary matrix endomorphism complex there
is then a degree-one cochain $h$ with
\[
                         z_a=d h+h d.
 \tag{F1R14.13}
\]
For a lower triangular entry, from $F_j$ to $F_i$ with
$i>j$, degree one means Dolbeault degree $1+j-i$.
Consequently every such entry with $i\ge j+2$ is zero.
The only possible lower entries are
\[
           f_j:=h_{j+1,j}\in
                    \mathcal A^{0,0}(\mathcal Hom(H_j,H_{j+1})).
 \tag{F1R14.14}
\]
The corresponding lower entries of $z_a$ are zero.
In that part of $d h+h d$, any term involving both
$\Delta$ and a further lower entry vanishes by the same
degree count. Hence $\bar\partial f_j=0$: these are actual
holomorphic bundle maps. The nearest upper entries
$\Delta_{j,j+1}$ are closed too, since the corresponding
entry of $d^2=0$ contains no product through an intermediate
cohomology degree. They represent actual classes in
$\operatorname{Ext}^2(H_{j+1},H_j)$.

Taking the diagonal entries of (F1R14.13) in cohomology
now gives only the adjacent products
\[
 [z_{a,jj}]=
       [\Delta_{j,j+1}f_j]+[f_{j-1}\Delta_{j-1,j}]
                         \quad\hbox{on }F_j.
 \tag{F1R14.15}
\]
This identity uses shifted morphisms, each displayed
factor of degree one; it fixes their composition signs.
The internal differential of $h_{jj}$ is a boundary.
All farther upper Postnikov entries remain part of the
actual model but cannot enter this diagonal equation,
because their required lower homotopy partners have
negative Dolbeault degree. This is the exact reason the
argument includes every actual higher coupling.

\paragraph{4. Full traces cancel separately at each actual slope.}
Decompose each $H_j$ by (F1R14.10), and write
$F_{j,s}=Q_{j,s}[-j]$ for the corresponding shifted block.
In a diagonal block of (F1R14.15), the contribution through
a different slope is zero in Ext by paragraph 2. Only the
same-slope components of $f_j$ and $\Delta_{j,j+1}$ remain.

The full Hochschild trace has the graded cyclic identity
\[
 \sigma_{P}^{p+q}(yx)
            =(-1)^{pq}\sigma_{Q}^{p+q}(xy)
 \quad(x:P\to Q[p],\ y:Q\to P[q]).
 \tag{F1R14.16}
\]
In the displayed notation composition includes its
required shift. To verify the formula, write the trace
as $\operatorname{Tr}((u\otimes1)\operatorname{AH})$,
move the natural degree-zero Atiyah--Hochschild character
past $x$, and then use graded cyclicity of the matrix
supertrace. Moving the two homogeneous maps past one
another supplies exactly $(-1)^{pq}$. This is the full
trace, not just its lowest component. The diagonal-kernel
naturality and this trace formula are the actual maps
of Buchweitz--Flenner \S5.1 and \S6.2.

Apply (F1R14.16) with $p=q=1$ to each adjacent same-slope
pair in (F1R14.15). Its two occurrences, on $F_{j,s}$
and $F_{j+1,s}$, cancel. Summing over all cohomological
degrees at a fixed slope gives
\[
 0=\sum_j\sigma^2_{F_{j,s}}
                    (\operatorname{ev}_{F_{j,s}}(a))
   =c^2_{w_s e^{\ell_s}}(a),\qquad
                 w_s=\sum_j(-1)^j r_{j,s}.
 \tag{F1R14.17}
\]
The shifted-character sign in $w_s$ is exactly the sign
of the perfect trace on $Q_{j,s}[-j]$. For two consecutive
degrees (F1R14.17) is equivalently the unshifted equality
$\sigma_A^2(\delta f)=\sigma_T^2(f[2]\delta)$ from the
single common map (F1R14.6).

\paragraph{5. The two-dimensional invariant trace kernel.}
The actual canonical triangle relations in $K_0(X)$ give
\[
                 [E]=\sum_j(-1)^j[H_j],\qquad
             \operatorname{ch}(E)=\sum_s w_s e^{\ell_s},
                         \qquad\sum_s w_s=N\ne0.
 \tag{F1R14.18}
\]
Choose a slope $s$ with $w_s\ne0$. What was just proved
is the exact kernel inclusion
\[
 \ker\left(\operatorname{ev}_E^2|_{K_E}\right)
       \subseteq
             \ker\left(c^2_{e^{\ell_s}}|_{K_E}\right).
 \tag{F1R14.19}
\]
R10's rank-four quotient calculation gives rank at least
two for the map on the right. Finite-dimensional rank
comparison in (F1R14.19) therefore gives
\[
                \operatorname{rank}
                    (\operatorname{ev}_E^2|_{K_E})\ge2.
 \tag{F1R14.20}
\]
Every element of this image has zero full trace, and
every specified coherent point linearization fixes it.
The identity-sector forgetful isomorphism of R05 identifies
these independent invariant ordinary Ext classes with
independent classes in the weak equivariant trace kernel.
This proves (F1R14.2) for the original object and lift.

\paragraph{Source, framework and scope record.}
The elementary new maps are the filtered forgetful fibre
(F1R14.4), its common Yoneda boundary, the actual Postnikov
matrix and the degree count (F1R14.12)--(F1R14.15), followed
by the slopewise cyclic trace. The diagonal-resolution roof
is the one explicitly constructed in the published
Buchweitz--Flenner Proposition 5.2.2; its proof and the
relevant trace formulas were reopened for this argument.
Mukai's actual isogeny-line construction is the published
input already used in R09. Mumford's original \S16 scan
was retrieved, rendered and read through both complete
relevant proofs and the signature corollary. Grieve's
published 2014 article, \S1.1, and Gulbrandsen's published
2008 Proposition 6.3 were read as additional primary
comparisons, not substituted for the retrieved original
index theorem.

GRR realizes the actual relations (F1R14.18). The isogeny
maps used in the vanishing proof retain their normalized
trace retractions and the Poincar\'e pullback identity of
F03. The point-action coherences are those of the original
lift; no new linear action or quotient object is inserted.
Adams operations act on the displayed index-polynomial
and character components by the codimension scalars of
F04. The final trace-kernel statement is proved using
the actual filtered object, beyond its $K_0$ realization.

This closes the arbitrary-coupling continuation of S-37
within the stated cohomology-filtration class. It does
not identify that class with all perfect objects, all
vector bundles, or every object in a triangulated category
generated by them. Objects with other actual cohomology
sheaves remain outside this theorem.
